\section{结论先行}
本报告重写后的核心结论是：光通信产业链必须按完整上下游理解，不能只看 AI 光模块。上游材料/基底决定光纤、激光芯片和调制器的物理边界；制造设备决定外延、硅光、主动耦合和测试良率；光芯片和光器件决定 1.6T/CPO 的战略稀缺性；光模块决定云厂商 AI capex 到收入利润的最直接弹性；光纤光缆、网络设备和下游应用决定运营商、DCI、FTTH 与企业网络的慢周期底盘。基于 2026-06-26 收盘后价格和 2026Q1 业绩折算，本报告覆盖 25 只可估值标的，另列产业链观察池。组合权重内在价值锚空间为 -47.1\%，市场共识调整后综合空间为 \textbf{-23.3\%}。

组合动作上，\textbf{中际旭创与新易盛}仍是最值得放在核心池里的两个光模块标的，原因是其 2026Q1 利润规模和毛利率已经证明海外 AI 订单不是纯主题；长飞光纤、亨通光电、中天科技、通鼎互联、永鼎股份提供光纤光缆和预制棒/线缆底盘，但应按慢周期现金流估值；中兴通讯和烽火通信映射网络设备/应用层，不能按纯光模块倍数追高；源杰科技、长光华芯、仕佳光子、光库科技、腾景科技代表上游光芯片/器件稀缺性；长芯博创、德科立、联特科技补足相干/高 beta 模块层；兆龙互连、意华股份、神宇股份补足高速线缆和连接器层，但需要客户认证、良率和 Q2/Q3 利润继续验证。

\begin{exhibitbox}[表：AStock 最终估值总表]
\centering
\tiny
\renewcommand{\arraystretch}{1.18}
\setlength{\tabcolsep}{2pt}
\begin{tabularx}{\exhibitboxwidth}{>{\bfseries\raggedright\arraybackslash}p{1.35cm} >{\raggedleft\arraybackslash}p{0.68cm} >{\raggedleft\arraybackslash}p{0.74cm} >{\raggedleft\arraybackslash}p{0.74cm} >{\raggedleft\arraybackslash}p{0.78cm} >{\raggedleft\arraybackslash}p{0.80cm} >{\raggedleft\arraybackslash}p{0.84cm} >{\centering\arraybackslash}p{0.78cm} >{\centering\arraybackslash}p{0.92cm} >{\raggedright\arraybackslash\sloppy\hspace{0pt}}X}
\toprule
\textbf{标的} & \textbf{现价} & \textbf{市值} & \textbf{26E EPS} & \textbf{27E EPS} & \textbf{综合目标} & \textbf{空间} & \textbf{动作} & \textbf{方法/证据} & \textbf{估值解释} \\
\midrule
\rowcolor{navy!6}中际旭创 300308 & 1253.89 & 13881 & 22.52 & 29.73 & 1229.56 & -1.9\% & \stance{riskamber}{中性} & \makecell[c]{PE/PEG\\A-} & 800G/1.6T 规模和盈利兑现最强，但当前估值已经反映海外 AI 需求持续高景气。 \\
\rowcolor{navy!6}新易盛 300502 & 566.00 & 5620 & 12.17 & 15.83 & 576.55 & +1.9\% & \stance{riskamber}{中性} & \makecell[c]{PE/PEG\\A-} & 高增长光模块弹性强，2025--2026 年交付较强，但估值对单一客户和 ASP 假设敏感。 \\
\rowcolor{accentblue!6}天孚通信 300394 & 318.00 & 2472 & 2.64 & 3.27 & 254.40 & -20.0\% & \stance{riskamber}{中性观察} & \makecell[c]{PE/PEG\\B+} & 高毛利精密器件平台价值明确，但增速相对模块龙头放缓。 \\
\rowcolor{accentblue!6}光迅科技 002281 & 240.34 & 1922 & 1.36 & 1.64 & 187.47 & -22.0\% & \stance{riskamber}{中性观察} & \makecell[c]{PE/PB/PS\\B} & 产品线和央企平台优势存在，但盈利能力显著低于 AI 光模块龙头。 \\
\rowcolor{riskamber!8}华工科技 000988 & 160.94 & 1606 & 2.56 & 3.02 & 128.75 & -20.0\% & \stance{riskamber}{中性观察} & \makecell[c]{PE/PB/PS\\B} & 混合业务提供盈利稳定性，但削弱纯 AI 光模块弹性。 \\
\rowcolor{riskamber!8}剑桥科技 603083 & 245.92 & 856 & 1.70 & 2.21 & 177.06 & -28.0\% & \stance{riskamber}{中性观察} & \makecell[c]{PE/PB/PS\\C+} & 反转弹性存在，但 Q1 利润分母仍偏小，难以支撑过高目标价。 \\
\rowcolor{riskamber!8}太辰光 300570 & 234.10 & 531 & 1.16 & 1.33 & 163.87 & -30.0\% & \stance{riskamber}{中性观察} & \makecell[c]{PE/PB/PS\\C} & 无源器件敞口真实，但 2026Q1 收入/利润下滑，估值需要交付修复。 \\
\rowcolor{riskamber!8}源杰科技 688498 & 1776.91 & 1526 & 9.09 & 12.54 & 1421.53 & -20.0\% & \stance{riskamber}{中性观察} & \makecell[c]{PE/PB/PS\\B} & 上游战略属性最强之一，但当前价格已经前置长期兑现路径。 \\
\rowcolor{accentblue!6}亨通光电 600487 & 110.81 & 2684 & 1.90 & 2.13 & 86.43 & -22.0\% & \stance{riskamber}{中性观察} & \makecell[c]{PE/PB/PS\\B} & 光纤光缆龙头兼具海缆和电网敞口，但 AI 数据中心纯度低于模块龙头。 \\
\rowcolor{accentblue!6}中天科技 600522 & 61.29 & 2071 & 1.13 & 1.27 & 47.81 & -22.0\% & \stance{riskamber}{中性观察} & \makecell[c]{PE/PB/PS\\B} & 电力海洋与通信线缆平台提供盈利韧性，但不能按纯 AI 光模块估值。 \\
\rowcolor{accentblue!6}长飞光纤 601869 & 543.00 & 4481 & 2.61 & 3.03 & 434.40 & -20.0\% & \stance{riskamber}{中性观察} & \makecell[c]{PE/PB/PS\\B+} & 预制棒/光纤/光缆一体化稀缺性较强，但数据中心需求需体现在利用率上。 \\
\rowcolor{accentblue!6}烽火通信 600498 & 76.22 & 975 & 0.14 & 0.16 & 53.35 & -30.0\% & \stance{riskamber}{中性观察} & \makecell[c]{PE/PB/PS\\B-} & 运营商光网络平台具备战略意义，但 Q1 利润分母薄，周期慢于云侧 AI capex。 \\
\rowcolor{accentblue!6}中兴通讯 000063 & 35.67 & 1731 & 1.23 & 1.37 & 31.06 & -12.9\% & \stance{riskamber}{中性} & \makecell[c]{PE/PB/PS\\A-} & 覆盖运营商、云网络和应用层，但光通信只是业务组合的一部分。 \\
\rowcolor{riskamber!8}仕佳光子 688313 & 179.70 & 812 & 1.12 & 1.40 & 125.79 & -30.0\% & \stance{riskamber}{中性观察} & \makecell[c]{PE/PB/PS\\B} & 无源光芯片与 WDM/AWG/硅光生态相关性强，但估值需要高利用率支撑。 \\
\rowcolor{riskamber!8}长光华芯 688048 & 396.00 & 698 & 0.13 & 0.18 & 277.20 & -30.0\% & \stance{riskamber}{中性观察} & \makecell[c]{PE/PB/PS\\C+} & 激光芯片期权高，但 Q1 利润分母仍小，基准情景要求快速认证和放量。 \\
\rowcolor{accentblue!6}光库科技 300620 & 362.43 & 903 & 0.78 & 0.98 & 116.95 & -67.7\% & \stance{riskred}{减持} & \makecell[c]{PE/PB/PS\\B} & 光器件和薄膜铌酸锂期权提高战略相关性，但股价已经计入较长兑现路径。 \\
\rowcolor{accentblue!6}通鼎互联 002491 & 36.15 & 444 & 0.18 & 0.20 & 26.03 & -28.0\% & \stance{riskamber}{中性观察} & \makecell[c]{PE/PB/PS\\B-} & 属于完整产业链底座，但盈利质量和 AI 数据中心纯度低于模块龙头。 \\
\rowcolor{accentblue!6}永鼎股份 600105 & 65.71 & 961 & 0.45 & 0.50 & 46.00 & -30.0\% & \stance{riskamber}{中性观察} & \makecell[c]{PE/PB/PS\\B} & 线缆和电信基础设施修复可见，但业务慢周期、纯度低于 AI 光模块。 \\
\rowcolor{accentblue!6}长芯博创 300548 & 262.65 & 759 & 1.88 & 2.34 & 189.11 & -28.0\% & \stance{riskamber}{中性观察} & \makecell[c]{PE/PB/PS\\B} & 光模块/器件平台具备数通敞口，但当前估值要求高速产品持续交付。 \\
\rowcolor{riskamber!8}德科立 688205 & 210.52 & 345 & 0.55 & 0.68 & 73.50 & -65.1\% & \stance{riskred}{减持} & \makecell[c]{PE/PB/PS\\B} & 补足 DCI 和运营商相干传输链条，但 Q1 利润基数仍偏小。 \\
\rowcolor{accentblue!6}腾景科技 688195 & 187.60 & 246 & 0.48 & 0.57 & 83.18 & -55.7\% & \stance{riskred}{减持} & \makecell[c]{PE/PB/PS\\B} & 精密光学与高速光系统相关，但规模小，估值需要持续订单转化。 \\
\rowcolor{riskamber!8}联特科技 301205 & 315.86 & 409 & 0.12 & 0.16 & 116.17 & -63.2\% & \stance{riskred}{减持} & \makecell[c]{PE/PB/PS\\C+} & 高速模块弹性大，但 Q1 利润分母很小，目标价高度依赖执行。 \\
\rowcolor{riskamber!8}兆龙互连 300913 & 50.75 & 174 & 0.59 & 0.70 & 22.44 & -55.8\% & \stance{riskred}{减持} & \makecell[c]{PE/PB/PS\\B} & 高速线缆/铜互连补足机柜内和企业网络层，但不是纯光模块替代品。 \\
\rowcolor{riskamber!8}意华股份 002897 & 84.50 & 162 & 0.95 & 1.10 & 49.92 & -40.9\% & \stance{riskred}{减持} & \makecell[c]{PE/PB/PS\\B-} & 连接器敞口属于数据中心互连链，但混合业务需要纯度折价。 \\
\rowcolor{riskamber!8}神宇股份 300563 & 24.73 & 48 & 0.32 & 0.37 & 12.95 & -47.7\% & \stance{riskred}{减持} & \makecell[c]{PE/PB/PS\\C+} & 线缆/互连拓宽产业链覆盖，但光通信纯度和 Q1 利润规模有限。 \\
\bottomrule
\end{tabularx}
\end{exhibitbox}
\sourcenote{行情：astock.quote\_service 2026-06-26 收盘后实时包；财务：astock.cli financials 2026Q1；估值模型：data/current\_valuation\_model\_20260626.json。}

\section{市场预期与券商对照先看结论}
炒股买的是预期，不只是过去一个季度。因此本报告在 AStock 自有目标价之外，额外给出基于 2026E 收入、2026E EPS、2027E 收入增长和成长性调整倍数的市场预期估值。该表回答“当前价格需要什么预期才能合理”，而不是重复财务历史。

\begin{exhibitbox}[表：市场预期估值桥]
\centering\tiny
\renewcommand{\arraystretch}{1.14}
\setlength{\tabcolsep}{2pt}
\begin{tabularx}{\exhibitboxwidth}{>{\bfseries\raggedright\arraybackslash}p{1.25cm} >{\raggedleft\arraybackslash}p{0.72cm} >{\raggedleft\arraybackslash}p{0.82cm} >{\raggedleft\arraybackslash}p{0.82cm} >{\raggedleft\arraybackslash}p{0.78cm} >{\raggedright\arraybackslash\sloppy\hspace{0pt}}p{1.65cm} >{\raggedleft\arraybackslash}p{0.82cm} >{\raggedleft\arraybackslash}p{0.82cm} >{\raggedright\arraybackslash\sloppy\hspace{0pt}}X}
\toprule
\textbf{标的} & \textbf{现价} & \textbf{2026E收入} & \textbf{收入增} & \textbf{2026E EPS} & \textbf{预期倍数} & \textbf{预期价} & \textbf{空间} & \textbf{预期驱动} \\
\midrule
中际旭创 300308 & 1253.89 & 847.67 & +32.0\% & 22.52 & PE 51.5x / PS 0.0x & 1159.42 & -7.5\% & 收入增长和毛利率兑现推动 EPS 预期，核心风险是海外云厂商订单持续性。 \\
新易盛 300502 & 566.00 & 362.52 & +30.0\% & 12.17 & PE 45.4x / PS 0.0x & 552.70 & -2.4\% & 收入增长和毛利率兑现推动 EPS 预期，核心风险是海外云厂商订单持续性。 \\
天孚通信 300394 & 318.00 & 55.43 & +24.0\% & 2.64 & PE 49.9x / PS 0.0x & 131.53 & -58.6\% & 收入增长和毛利率兑现推动 EPS 预期，核心风险是海外云厂商订单持续性。 \\
光迅科技 002281 & 240.34 & 126.06 & +20.0\% & 1.36 & PE 38.2x / PS 4.0x & 54.27 & -77.4\% & 市场预期来自收入增长、毛利率改善和业务纯度提升。 \\
华工科技 000988 & 160.94 & 170.63 & +18.0\% & 2.56 & PE 37.8x / PS 3.9x & 86.03 & -46.5\% & 市场预期来自收入增长、毛利率改善和业务纯度提升。 \\
剑桥科技 603083 & 245.92 & 64.33 & +30.0\% & 1.70 & PE 51.1x / PS 6.0x & 97.08 & -60.5\% & 市场预期主要来自高增长和战略稀缺性，需用客户认证和收入规模验证。 \\
太辰光 300570 & 234.10 & 13.63 & +15.0\% & 1.16 & PE 37.4x / PS 0.9x & 25.19 & -89.2\% & 收入增长预期主要来自运营商/云侧线缆项目，估值上限受现金流和线缆价格约束。 \\
源杰科技 688498 & 1776.91 & 15.45 & +38.0\% & 9.09 & PE 76.1x / PS 6.2x & 375.41 & -78.9\% & 市场预期主要来自高增长和战略稀缺性，需用客户认证和收入规模验证。 \\
亨通光电 600487 & 110.81 & 741.27 & +12.0\% & 1.90 & PE 25.3x / PS 0.9x & 35.55 & -67.9\% & 收入增长预期主要来自运营商/云侧线缆项目，估值上限受现金流和线缆价格约束。 \\
中天科技 600522 & 61.29 & 547.59 & +12.0\% & 1.13 & PE 25.3x / PS 0.9x & 22.59 & -63.1\% & 收入增长预期主要来自运营商/云侧线缆项目，估值上限受现金流和线缆价格约束。 \\
\bottomrule
\end{tabularx}
\end{exhibitbox}
\sourcenote{完整 25 只标的预期估值见第 8 章和 data/market\_expectation\_valuation\_20260626.json。}

市场情绪同样是估值方法论的一部分。市场可以错，但当前价、成交额、券商目标和共识定价是可观测信息。本报告把它们作为市场隐含预期锚，和内在价值锚、券商锚共同形成综合目标价。若内在价值锚远低于股价但成交和共识仍强，动作不会机械写成减持，而会标注为中性观察并给出情绪溢价失效条件。

\begin{exhibitbox}[表：市场隐含预期与情绪锚]
\centering\tiny
\renewcommand{\arraystretch}{1.13}
\setlength{\tabcolsep}{2pt}
\begin{tabularx}{\exhibitboxwidth}{>{\bfseries\raggedright\arraybackslash}p{1.25cm} >{\raggedleft\arraybackslash}p{0.72cm} >{\raggedleft\arraybackslash}p{0.78cm} >{\raggedright\arraybackslash\sloppy\hspace{0pt}}p{1.45cm} >{\raggedleft\arraybackslash}p{0.72cm} >{\centering\arraybackslash}p{0.78cm} >{\raggedleft\arraybackslash}p{0.78cm} >{\raggedleft\arraybackslash}p{0.78cm} >{\raggedright\arraybackslash\sloppy\hspace{0pt}}p{1.25cm} >{\raggedleft\arraybackslash}p{0.82cm} >{\raggedright\arraybackslash\sloppy\hspace{0pt}}X}
\toprule
\textbf{标的} & \textbf{内在锚} & \textbf{市场锚} & \textbf{当前隐含倍数} & \textbf{成交分位} & \textbf{情绪} & \textbf{券商锚} & \textbf{综合目标} & \textbf{权重} & \textbf{空间} & \textbf{动作逻辑} \\
\midrule
中际旭创 300308 & 1337.79 & 1003.11 & PE 55.7倍 / PS 16.4倍 & +100.0\% & 强共识/高拥挤 & 1062.50 & 1229.56 & 内在65\% / 市场20\% / 券商15\% & -1.9\% & 动作由综合目标价相对当前价决定。 \\
新易盛 300502 & 633.04 & 452.80 & PE 46.5倍 / PS 15.5倍 & +96.0\% & 强共识/高拥挤 & 496.76 & 576.55 & 内在65\% / 市场20\% / 券商15\% & +1.9\% & 动作由综合目标价相对当前价决定。 \\
天孚通信 300394 & 147.20 & 254.40 & PE 120.5倍 / PS 44.6倍 & +84.0\% & 强共识/高拥挤 & 286.81 & 254.40 & 内在57\% / 市场28\% / 券商15\% & -20.0\% & 市场高成交已给出情绪溢价，维持中性观察并等待利润/现金流追认。 \\
光迅科技 002281 & 58.46 & 187.47 & PE 176.2倍 / PS 15.2倍 & +80.0\% & 强共识/高拥挤 & 123.86 & 187.47 & 内在47\% / 市场38\% / 券商15\% & -22.0\% & 市场高成交已给出情绪溢价，维持中性观察并等待利润/现金流追认。 \\
华工科技 000988 & 92.79 & 128.75 & PE 62.9倍 / PS 9.4倍 & +76.0\% & 强共识/高拥挤 & 133.99 & 128.75 & 内在47\% / 市场38\% / 券商15\% & -20.0\% & 市场高成交已给出情绪溢价，维持中性观察并等待利润/现金流追认。 \\
剑桥科技 603083 & 105.68 & 177.06 & PE 144.7倍 / PS 13.3倍 & +40.0\% & 活跃共识 & 165.68 & 177.06 & 内在47\% / 市场38\% / 券商15\% & -28.0\% & 市场高成交已给出情绪溢价，维持中性观察并等待利润/现金流追认。 \\
太辰光 300570 & 26.60 & 163.87 & PE 201.8倍 / PS 39.0倍 & +52.0\% & 活跃共识 & 148.32 & 163.87 & 内在38\% / 市场47\% / 券商15\% & -30.0\% & 市场高成交已给出情绪溢价，维持中性观察并等待利润/现金流追认。 \\
源杰科技 688498 & 434.08 & 1421.53 & PE 195.5倍 / PS 98.8倍 & +64.0\% & 强共识/高拥挤 & 1164.00 & 1421.53 & 内在47\% / 市场38\% / 券商15\% & -20.0\% & 市场高成交已给出情绪溢价，维持中性观察并等待利润/现金流追认。 \\
亨通光电 600487 & 36.95 & 86.43 & PE 58.3倍 / PS 3.6倍 & +88.0\% & 强共识/高拥挤 & 97.88 & 86.43 & 内在38\% / 市场47\% / 券商15\% & -22.0\% & 市场高成交已给出情绪溢价，维持中性观察并等待利润/现金流追认。 \\
中天科技 600522 & 23.35 & 47.81 & PE 54.1倍 / PS 3.8倍 & +92.0\% & 强共识/高拥挤 & 49.61 & 47.81 & 内在38\% / 市场47\% / 券商15\% & -22.0\% & 市场高成交已给出情绪溢价，维持中性观察并等待利润/现金流追认。 \\
\bottomrule
\end{tabularx}
\end{exhibitbox}
\sourcenote{完整市场情绪锚见第 8 章和 data/market\_sentiment\_anchor\_20260626.json。}

公开券商/一致预期对照用于回答“市场怎么想”。本报告只使用可追溯的公开一致预期、券商摘要或原始入口；没有披露 2026E/2027E 收入、利润、EPS 或目标价的字段，统一标记为未披露，不用 AStock 假设回填。

\begin{exhibitbox}[表：公开券商/一致预期对照]
\centering\tiny
\renewcommand{\arraystretch}{1.14}
\setlength{\tabcolsep}{2pt}
\begin{tabularx}{\exhibitboxwidth}{>{\bfseries\raggedright\arraybackslash}p{1.25cm} >{\raggedright\arraybackslash\sloppy\hspace{0pt}}p{1.75cm} >{\centering\arraybackslash}p{0.72cm} >{\centering\arraybackslash}p{0.52cm} >{\raggedleft\arraybackslash}p{0.75cm} >{\raggedleft\arraybackslash}p{0.75cm} >{\raggedleft\arraybackslash}p{0.75cm} >{\raggedleft\arraybackslash}p{0.75cm} >{\raggedright\arraybackslash\sloppy\hspace{0pt}}X}
\toprule
\textbf{标的} & \textbf{来源} & \textbf{评级} & \textbf{证据} & \textbf{券商均值} & \textbf{高值} & \textbf{低值} & \textbf{券商空间} & \textbf{预测说明} \\
\midrule
中际旭创 300308 & 英为财情一致预期/21财经券商聚合 & 偏积极 & B & 1062.50 & 1650.00 & 430.00 & -15.3\% & 公开页面披露目标价区间；部分券商聚合提到华泰证券预计 2026E 归母净利润约 284.7 亿元。 \\
新易盛 300502 & 英为财情一致预期 & 偏积极 & B & 496.76 & 701.00 & 264.29 & -12.2\% & 公开页面披露目标价区间，未完整披露 2026E 收入/利润明细。 \\
天孚通信 300394 & 英为财情一致预期/同花顺预测摘要 & 偏积极 & B & 286.81 & 419.00 & 115.71 & -9.8\% & 公开页面披露目标价；同花顺摘要披露券商未来几年净利润预测但非完整原文。 \\
光迅科技 002281 & 英为财情一致预期 & 偏积极 & B & 123.86 & 166.00 & 78.31 & -48.5\% & 公开页面披露目标价区间，未完整披露 2026E 收入/利润明细。 \\
华工科技 000988 & 英为财情一致预期 & 偏积极 & B- & 133.99 & 157.98 & 110.00 & -16.7\% & 公开页面披露目标价区间，样本较少。 \\
剑桥科技 603083 & 英为财情一致预期 & 偏积极 & C+ & 165.68 & 165.68 & 165.68 & -32.6\% & 公开页面披露单一目标价，样本很少。 \\
太辰光 300570 & 英为财情一致预期 & 样本有限 & C+ & 148.32 & 154.96 & 138.00 & -36.6\% & 公开检索显示目标价样本有限，需以原始研报复核。 \\
源杰科技 688498 & 英为财情一致预期 & 偏积极 & B- & 1164.00 & 1253.39 & 1097.93 & -34.5\% & 公开页面披露目标价区间，未完整披露 2026E 收入/利润明细。 \\
亨通光电 600487 & 英为财情一致预期 & 偏积极 & B & 97.88 & 132.00 & 77.33 & -11.7\% & 公开页面披露目标价区间，未完整披露 2026E 收入/利润明细。 \\
中天科技 600522 & 英为财情一致预期 & 偏积极 & B & 49.61 & 66.00 & 26.40 & -19.1\% & 公开页面披露目标价区间，未完整披露 2026E 收入/利润明细。 \\
\bottomrule
\end{tabularx}
\end{exhibitbox}
\sourcenote{完整券商/一致预期对照见第 7 章和 data/broker\_consensus\_snapshot\_20260626.json。券商目标价不直接作为 AStock 目标价。}

\section{评级与权重方法}
本报告用五类动作标签约束组合行为：买入表示综合目标价较当前价有 20\% 以上空间且证据质量不低于 B；增持表示空间 8--20\%；中性表示空间介于 -15\% 至 8\%；中性观察表示内在价值锚偏低但市场成交、共识定价或情绪溢价仍强，需要等待利润和现金流追认；减持表示综合空间低于 -15\% 且情绪支撑不足或正在破裂。权重不是建议仓位，而是用于计算产业链组合赔率的研究权重：光模块与器件 58\%，上游光芯片/光源 13\%，光纤光缆 17\%，网络设备与应用 12\%。设备和材料中缺少稳定 PE 分母的对象进入观察池，不参与目标价加权。

\begin{exhibitbox}[表：组合执行摘要 / Portfolio Action Summary]
\centering
\small
\begin{tabularx}{\exhibitboxwidth}{>{\bfseries\raggedright\arraybackslash}p{2.5cm} >{\raggedright\arraybackslash\sloppy\hspace{0pt}}X >{\raggedright\arraybackslash}p{3.2cm}}
\toprule
\textbf{动作} & \textbf{标的} & \textbf{执行含义} \\
\midrule
核心持有/增持观察 & 中际旭创、新易盛 & 只在 Q2 订单、毛利率与现金流继续验证时提高权重 \\
慢周期底盘 & 长飞光纤、亨通光电、中天科技、通鼎互联、永鼎股份、中兴通讯、烽火通信 & 用现金流、运营商/海缆/FTTH 项目节奏和网络设备利润率校验 \\
上游与相干期权 & 源杰科技、长光华芯、仕佳光子、光库科技、腾景科技、长芯博创、德科立、联特科技 & 只在认证、良率、收入占比和毛利率同时改善时提高估值容忍度 \\
线缆/连接器观察 & 兆龙互连、意华股份、神宇股份、光迅科技、华工科技、剑桥科技、太辰光 & 当前价格需要更多利润分母或业务纯度证明 \\
\bottomrule
\end{tabularx}
\end{exhibitbox}

\section{下一季度最重要的三个问题}
第一，中际旭创(300308), 新易盛(300502), 天孚通信(300394) 的 2026Q2 净利润是否至少达到 Q1 的 1.05 倍，这是判断 800G 订单是否继续外溢的最低门槛。第二，1.6T 资格认证是否转化为批量订单，而不是停留在样品和送样阶段。第三，上游光芯片/器件、光纤光缆和网络设备是否出现独立验证，而不是全部依赖光模块情绪扩散；如果只有模块端景气、其他层级没有收入和现金流同步，完整产业链的估值就应继续分层。
